\documentclass[11pt,a4paper]{article}

\usepackage[margin=1.05in]{geometry}
\usepackage{amsmath,amssymb,amsthm}
\usepackage{mathtools}
\usepackage{stmaryrd}
\usepackage{microtype}
\setlength{\emergencystretch}{2.5em}
\usepackage{enumitem}
\usepackage{booktabs}
\usepackage{array}
\usepackage{tikz-cd}
\usepackage{xcolor}
\usepackage[colorlinks=true,linkcolor=blue!45!black,urlcolor=blue!60!black,citecolor=blue!45!black]{hyperref}

\newtheorem{theorem}{Theorem}[section]
\newtheorem{proposition}[theorem]{Proposition}
\newtheorem{lemma}[theorem]{Lemma}
\newtheorem{corollary}[theorem]{Corollary}
\newtheorem{conjecture}[theorem]{Conjecture}
\theoremstyle{definition}
\newtheorem{definition}[theorem]{Definition}
\newtheorem{example}[theorem]{Example}
\theoremstyle{remark}
\newtheorem{remark}[theorem]{Remark}
\newtheorem{warning}[theorem]{Warning}
\newtheorem{question}[theorem]{Question}

\newcommand{\FF}{\mathbb{F}}
\newcommand{\NN}{\mathbb{N}}
\newcommand{\ZZ}{\mathbb{Z}}
\newcommand{\QQ}{\mathbb{Q}}
\newcommand{\RR}{\mathbb{R}}
\newcommand{\cat}[1]{\mathbf{#1}}
\newcommand{\Set}{\cat{Set}}
\newcommand{\FinSet}{\cat{FinSet}}
\newcommand{\Poly}{\cat{Poly}}
\newcommand{\Cont}{\cat{Cont}}
\newcommand{\Spc}{\cat{Spc}}
\newcommand{\iso}{\mathrm{iso}}
\newcommand{\con}{\mathrel{\vartriangleright}}
\newcommand{\norm}[1]{\lvert #1 \rvert}
\newcommand{\dd}{\,d}
\newcommand{\Had}{\odot}
\DeclareMathOperator{\GL}{GL}
\DeclareMathOperator{\Aut}{Aut}
\DeclareMathOperator{\List}{List}
\DeclareMathOperator{\Bag}{Bag}
\newcommand{\qbin}[2]{\genfrac{[}{]}{0pt}{}{#1}{#2}_{q}}

\definecolor{okgreen}{rgb}{0.0,0.45,0.15}
\definecolor{nored}{rgb}{0.65,0.0,0.0}
\newcommand{\OK}{\textcolor{okgreen}{\textbf{verified}}}
\newcommand{\NO}{\textcolor{nored}{\textbf{fails}}}

\title{\bfseries Hypergeometric Series, Containerwise\\[0.4em]
\large The missing column: superposition rather than convolution,
and why the $q$-side is the honest one}
\author{}
\date{}

\begin{document}
\maketitle

\begin{abstract}
Hypergeometric series look alien to generating-function intuition, and this note
argues that the reason is structural rather than notational: they do not live in the
monoidal structure that \texttt{q-dirichlet.tex} tabulates. Four claims. (i) The
obligatory $n!$ in ${}_pF_q$ is the $S_n$-quotient --- the notation is native to
$\Bag$-world, which is why the geometric series sits outside it. (ii) The
\emph{upstairs} parameters are containers, $1/(1-x)^a = \List^a$, but they are
combined by the \emph{Hadamard} product --- superposition of structures on one set
--- which is the categorical product in the species category and not a Day
convolution at all. That is a missing \emph{column} of the table where
\texttt{list-and-bag.tex} supplied a missing row. (iii) The \emph{downstairs}
parameters are containerish only at $b = 1$, where $(1)_n = \norm{S_n}$; this is
exactly why ${}_0F_{s-1}(\,;1,\dots,1;z)$ was the family for which
\texttt{list-and-bag.tex} Thm.~5.1 was clean, and predicts that no such theorem
exists for ${}_2F_1$. (iv) The defining property, rationality of $c_{n+1}/c_n$, is a
level-by-level bijection once denominators are cleared. The comparison with Fomin's
local rules (\texttt{lascoux/}) is \emph{narrower} than it first looks --- Fomin's
rule is vertical and two-dimensional, a bare recurrence is horizontal and
one-dimensional --- and \S5 states the comparison that survives: both are
commutator identities categorified on a two-dimensional growth diagram, with the
Wilf--Zeilberger pair supplying the second dimension. Finally: at $q$ the obligatory $(q;q)_n$ \emph{is}
$\norm{\GL_n(\FF_q)}$ up to the torus factor, so basic hypergeometric series are more
container-native than classical ones, and $q$-Vandermonde is already the Levi row of
the decategorification table.
\end{abstract}

\section{What the notation is hiding}

\begin{definition}[the notation]\label{def:pFq}
With $(b)_n = b(b+1)\cdots(b+n-1) = \Gamma(b+n)/\Gamma(b)$ the rising factorial,
\[
  {}_pF_q\!\left(\begin{matrix} a_1,\dots,a_p \\ b_1,\dots,b_q\end{matrix}; z\right)
  \;=\; \sum_{n \ge 0} c_n z^n,
  \qquad
  c_n \;=\; \frac{(a_1)_n \cdots (a_p)_n}{(b_1)_n\cdots(b_q)_n}\cdot\frac{1}{n!} .
\]
\end{definition}

Three features of $c_n$ are usually presented as conventions. They are not, and the
rest of this note is an account of what each one is.

\begin{enumerate}[nosep,label=(\alph*)]
  \item There is always a spare $1/n!$, whatever $p$ and $q$ are. \hfill(\S2)
  \item The $a_i$ and $b_j$ are evaluated at the \emph{same} $n$, and multiplied.
        \hfill(\S3)
  \item The $b_j$ appear in a denominator, with no companion construction upstairs.
        \hfill(\S4)
\end{enumerate}

\section{The spare $n!$ is the $S_n$}

\begin{proposition}[\emph{ours}, immediate]\label{prop:native}
Definition~\ref{def:pFq} is a statement about exponential generating functions. Write
$c_n = A_n/n!$ with $A_n = \prod_i (a_i)_n / \prod_j (b_j)_n$; then ${}_pF_q$ is the
EGF of the sequence $(A_n)$. Hence hypergeometric notation is native to
$\Bag$-world: the $n!$ is the $\norm{S_n}$ that a quotient container divides by, taken
once, unconditionally.
\end{proposition}

This is the whole of (a). The $n!$ is not an extra factor bolted onto a product of
Pochhammers --- it is the normalisation that says ``now read the result as a species
cardinality''. Two consequences worth stating, because both look like notational
quirks and are not:

\begin{example}[the exponential is the emptiest one]\label{ex:exp}
${}_0F_0(\,;\,;z) = \sum z^n/n! = e^z$. With no parameters at all, what is left is
exactly the $S_n$-quotient. In the language of \texttt{list-and-bag.tex}, ${}_0F_0$
\emph{is} $\Bag$.
\end{example}

\begin{example}[the geometric series is outside the notation]\label{ex:geo}
$\sum z^n = 1/(1-z)$ is not any ${}_0F_{-1}$: the notation begins at $q = 0$. It is
${}_1F_0(1;\,;z)$, with $(1)_n = n!$ upstairs cancelling the obligatory $n!$
downstairs. So $\List$ can be written only by \emph{undoing} the quotient with a
parameter, and the asymmetry between Examples~\ref{ex:exp} and \ref{ex:geo} is the
$\List/\Bag$ asymmetry in notational form.
\end{example}

\begin{remark}\label{rem:struggle}
If hypergeometric notation feels like it is doing something arbitrary, this is the
first thing to fix. It is not a notation for power series. It is a notation for
\emph{species}, in which the ``divide by $n!$'' step has been performed in advance and
made invisible.
\end{remark}

\section{Upstairs: containers, combined by superposition}

\subsection{A single parameter is a power of $\List$}

\begin{proposition}\label{prop:pochlist}
For $a \in \NN$ let $L^a$ be the species of $a$-tuples of linear orders: $L^a[n]$ is
the set of ways to distribute $[n]$ into an ordered $a$-tuple of lists. Then
\[
  \norm{L^a[n]} = (a)_n,
  \qquad\text{equivalently}\qquad
  \sum_{n\ge0}(a)_n\frac{x^n}{n!} = \frac{1}{(1-x)^a} .
\]
As a container, $L^a = \List^{\times a}$, the $a$-fold product of $\List$ in the sense
of \texttt{q-container.tex} Def.~(product of containers).
\end{proposition}

\begin{proof}
$(1-x)^{-a} = \sum_n \binom{a+n-1}{n}x^n = \sum_n (a)_n x^n/n!$, and the species
product decategorifies to multiplication of EGFs, so the EGF of $L^a$ is
$(1/(1-x))^a$.
\end{proof}

\begin{example}\label{ex:a2n2}
$a = 2$, $n = 2$: $(2)_2 = 2\cdot3 = 6$, and the pairs of lists on $\{1,2\}$ are
$12\mid\varnothing$, $21\mid\varnothing$, $\varnothing\mid12$, $\varnothing\mid21$,
$1\mid2$, $2\mid1$. \OK
\end{example}

So a rising factorial upstairs is a perfectly good container, and for integer $a$ the
hypergeometric parameter is a \emph{dimension}: how many lists.

\subsection{Several parameters: the Hadamard product}

Feature (b) is the interesting one. In $c_n$ the factors $(a_1)_n,\dots,(a_p)_n$ are
evaluated at one and the same $n$ and multiplied. That is \emph{not} the Cauchy
product.

\begin{definition}[Cartesian/Hadamard product of species]\label{def:had}
For species $F, G$ let $(F \Had G)[n] = F[n] \times G[n]$, with the diagonal
$S_n$-action. This is the categorical product in the functor category
$[\FinSet^{\iso}, \Set]$, i.e.\ the \emph{pointwise} product of presheaves. Its effect
on EGF coefficients is $A_n \mapsto A_n B_n$ with one $n!$ retained.
\end{definition}

\begin{proposition}[\emph{ours}]\label{prop:had-upstairs}
The numerator of $c_n$ in Definition~\ref{def:pFq} is the EGF coefficient sequence of
\[
  L^{a_1} \Had L^{a_2} \Had \cdots \Had L^{a_p},
\]
and $c_n$ is that sequence read as an EGF. Hypergeometric notation is therefore:
\emph{Hadamard-multiply powers of $\List$, then read the result in $\Bag$-world.}
\end{proposition}

Containerwise the Hadamard product is a pullback. Let $C = (S \con P)$ and
$D = (T \con Q)$ be containers and let $\norm{\cdot} : S \to \NN$, $s \mapsto \norm{P s}$,
be the arity map.

\begin{definition}[Hadamard product of containers]\label{def:had-cont}
$C \Had D = \bigl(S \times_{\NN} T \con (s,t)\mapsto P s\bigr)$, the fibre product of
the arity maps, with positions inherited from either factor (they agree in
cardinality by construction).
\end{definition}

\begin{remark}[the slogan]\label{rem:slogan}
Cauchy product: \emph{split} the positions, one factor on each part. Hadamard
product: \emph{share} the positions, both factors on all of them. Superposition, not
concatenation.
\end{remark}

\subsection{Why this is a missing column and not a missing row}

\begin{warning}[\emph{the central claim of this note}]\label{warn:column}
The Hadamard product is not a Day convolution. Day convolution along
$(\mathcal{A},\otimes)$ is the unique monoidal structure making Yoneda strong
monoidal, and it uses $\otimes$ on objects of $\mathcal{A}$; the table of
\texttt{q-dirichlet.tex} \S2 varies $\mathcal{A}$ and reads off $\sqcup$ (species
product, EGF), $\times$ on $\NN_{>0}$ (Dirichlet), $\times$ on $\FinSet^{\iso}$
(arithmetic product). Definition~\ref{def:had} is none of these: it is the
\emph{categorical product} in the presheaf category, which exists for formal reasons
and is not a convolution along anything.
\end{warning}

\begin{center}
\small
\renewcommand{\arraystretch}{1.2}
\begin{tabular}{@{}lll@{}}
\toprule
monoidal structure on species & on positions & on coefficients \\
\midrule
Cauchy / Day along $\sqcup$ & split & $\sum_{i+j=n}\binom{n}{i}A_iB_j$ \\
Dirichlet / Day along $\times$ & multiply & $\sum_{ij=n}A_iB_j$ \\
\textbf{Hadamard / pointwise product} & \textbf{share} & $\boxed{A_nB_n}$ \\
\bottomrule
\end{tabular}
\end{center}

\begin{corollary}[\emph{ours} --- the answer to ``why do these feel alien''] \label{cor:alien}
Generating-function intuition trained on the Day/Cauchy structure reads
``multiply the series'' as ``concatenate the structures''. Hypergeometric parameters
multiply in the other monoidal structure entirely, in which multiplying means
overlaying two structures on one set of positions. The notation does not record which
structure is in force, and so gives no clue that the usual intuition is inapplicable.
\end{corollary}

\begin{question}\label{q:had-q}
What is the Hadamard product at $q$? Definition~\ref{def:had-cont} needs an arity map
to pull back along; for $q$-containers the natural one is
$\norm{\cdot} : S \to \NN$, $s \mapsto \dim P s$, and the fibre product is over
dimension. Is the class of $q$-containers closed under it --- and does it
decategorify to the coefficientwise product of $q$-series? This is cheap to check and
is the first thing to do with this note.
\end{question}

\section{Downstairs: only $b = 1$ is a quotient}

Feature (c). Dividing by $(b)_n$ coefficientwise is division in the Hadamard ring. A
quotient \emph{container}, by contrast, requires a group that acts on positions, and
then divides by its order.

\begin{proposition}[\emph{ours}]\label{prop:downstairs}
\begin{enumerate}[nosep,label=(\roman*)]
  \item $b = 1$: $(1)_n = n! = \norm{S_n}$, and $S_n$ acts on the $n$ positions.
        Division by $(1)_n$ is the formation of a quotient container in Kock's sense
        \cite{Kock} --- a polynomial functor over groupoids --- and is structural.
  \item $b \ge 2$: $(b)_n = \norm{L^b[n]}$ is a cardinality, but no group of that
        order acts canonically on anything in sight. Division by $(b)_n$ is
        arithmetic in the coefficient ring, not a quotient.
  \item $b \notin \NN$: $(b)_n = \Gamma(b+n)/\Gamma(b)$ is not a cardinality at all.
\end{enumerate}
\end{proposition}

\begin{corollary}[why \texttt{list-and-bag.tex} Thm.~5.1 was clean]\label{cor:whyclean}
The family $F_s(z) = \sum z^n/(n!)^s = {}_0F_{s-1}(\,;1,\dots,1;z)$ is exactly the
hypergeometric family whose \emph{entire} denominator consists of case (i): every
division it performs is an honest $S_n$-quotient. That is why the exponent in
$\Gamma(\tfrac12)^{2-s}$ counts something --- it counts quotients taken --- and the
statement has a container reading at all.
\end{corollary}

\begin{conjecture}[\emph{ours}, a prediction rather than a hope]\label{conj:no2F1}
There is no analogue of \texttt{list-and-bag.tex} Thm.~5.1 for ${}_2F_1$ in which the
answer is a power of a single $\Gamma$-value and the exponent counts a container
construction. The Mellin transform of ${}_2F_1(a,b;c;-x^2)$ exists and is a ratio of
four $\Gamma$'s (Barnes), but by Proposition~\ref{prop:downstairs}(ii) the
$c$-dependence is not quotient data, so no exponent can be counting quotients.
\end{conjecture}

The point of recording Conjecture~\ref{conj:no2F1} is that it is the honest limit of
the programme. The container reading explains why one family worked; it should also
say which families cannot.

\section{Local rules: the comparison, and its limits}

\subsection{The recurrence as a bijection}

\begin{definition}\label{def:hyp}
A series $\sum c_n z^n$ is \emph{hypergeometric} when $c_{n+1}/c_n$ is a rational
function of $n$. Definition~\ref{def:pFq} is the general solution:
$c_{n+1}/c_n = \prod_i(n+a_i)\big/\bigl(\prod_j(n+b_j)\,(n+1)\bigr)$.
\end{definition}

Clear denominators. Writing $c_{n+1}/c_n = P(n)/Q(n)$,
\begin{equation}\label{eq:localrule}
  Q(n)\cdot c_{n+1} \;=\; P(n)\cdot c_n .
\end{equation}

\begin{proposition}[\emph{ours}]\label{prop:localrule}
When $P$ and $Q$ take non-negative integer values on $\NN$ --- which happens exactly
when all $a_i, b_j \in \NN$, the case of \S3--\S4 --- \eqref{eq:localrule} is a
bijection statement: $Q(n)$ copies of the level-$(n+1)$ shapes are in bijection with
$P(n)$ copies of the level-$n$ shapes.
\end{proposition}

It is tempting to stop here and say that a hypergeometric series is therefore a
container presented by a local rule in the sense of \texttt{lascoux/}. That is what an
earlier draft of this note said, and it is too strong. The next two subsections say
what Fomin's rule actually is and then what survives of the comparison.

\subsection{What Fomin's local rule is}

A growth diagram is a grid whose lattice \emph{points} carry partitions, adjacent
points differing by at most one box, and whose \emph{faces} carry a label
$x \in \{0,1\}$ --- the dots of a permutation matrix. Three corners of a face plus the
label determine the fourth:
\[
\begin{array}{ccc}
 \rho & \!\!\text{\textemdash\textemdash}\!\! & \nu \\[-0.2em]
 \mid & x & \mid \\[-0.2em]
 \mu & \!\!\text{\textemdash\textemdash}\!\! & \lambda
\end{array}
\]
Fill the left and top boundary with empty partitions, iterate over all faces, and the
right and bottom boundaries carry the RSK tableaux of the permutation. The six cases,
as formalised in \texttt{lascoux/lean/Fomin/LocalRule.lean} (thesis \S2.1.6):

\begin{center}
\small
\renewcommand{\arraystretch}{1.15}
\begin{tabular}{@{}clll@{}}
\toprule
 & condition & $\lambda$ & \\
\midrule
1 & $\mu=\nu=\rho$, $x=1$ & $\rho$ + box in row $0$ & the dot \\
2 & $\mu=\nu=\rho$, $x=0$ & $\rho$ & nothing \\
3 & $\mu=\rho$, $\nu\ne\rho$ & $\nu$ & shadow passes vertically \\
4 & $\nu=\rho$, $\mu\ne\rho$ & $\mu$ & shadow passes horizontally \\
5 & $\mu,\nu\ne\rho$, $\mu\ne\nu$ & $\mu \sqcup \nu$ & two shadows cross \\
6 & $\mu,\nu\ne\rho$, $\mu=\nu$ & $\mu$ + box in row $\mathrm{diffRow}(\mu,\rho)+1$ & the bounce \\
\bottomrule
\end{tabular}
\end{center}

\begin{remark}[all the content is in rule 6]\label{rem:bounce}
Five of the six cases are free: two identities, two projections and a lattice join,
needing nothing about Young diagrams beyond the lattice structure. Rule 6 is the only
one that computes --- it must find the row the new box lands in --- and it is where
Schensted bumping lives. The Lean file exhibits this asymmetry structurally:
\texttt{diffRow} and the lemma that the bounce always has somewhere to land
(\texttt{isAddableRow\_succ\_of\_addBox}) exist for rule 6 alone, while rule 5 needs
only cardinality arithmetic on a join. In one sentence: \emph{the unbounded bumping
cascade of Schensted insertion compresses into a single row index}, and the
cascade is recovered by iteration over the grid rather than by recursion inside
the rule.
\end{remark}

\begin{remark}[why \emph{vertical}]\label{rem:vertical}
The rule has type $\ZZ_{\ge0}\times(D\alpha \cap D\beta) \to U\alpha \cap U\beta$. A
container morphism $(S\con P) \to (T \con Q)$ is a pair: $f : S \to T$ on shapes and,
\emph{contravariantly}, $g : \prod_{s} Q(f s) \to P s$ on positions. \emph{Vertical}
means $f = \mathrm{id}$: the shapes are untouched and only the fibres move. Fomin's
rule is vertical, and that is the precise content of the word \emph{local} --- the
rule is uniform in $(\alpha,\beta)$, one formula with six case distinctions and no
further dependence on which partitions one is at.
\end{remark}

\begin{remark}[what it decategorifies to]\label{rem:weyl}
Counting the fibres of a vertical isomorphism gives
$\norm{U\alpha \cap U\beta} = \delta_{\alpha\beta} + \norm{D\alpha\cap D\beta}$, which
coefficientwise in $\ZZ[\text{Young}]$ is
\[
  DU - UD = I,
\]
Stanley's differential-poset axiom, i.e.\ the Weyl commutation relation. The
$\delta_{\alpha\beta}$ is the label $x=1$ being legal at exactly one face shape. So
Fomin's local rule is a \emph{bijection witnessing a commutator identity}, and the
step-for-step correspondence between applications of the rule and applications of the
commutation relation (thesis \S4.3.1) is the statement that the algebraic proof is the
decategorification of the diagrammatic one.
\end{remark}

\subsection{The comparison that survives}

\begin{warning}[the analogy is narrower than Proposition~\ref{prop:localrule} suggests]\label{warn:analogy}
Two mismatches, both structural.
\begin{enumerate}[nosep,label=(\roman*)]
  \item \textbf{Vertical versus horizontal.} Fomin's rule fixes the shapes and
        bijects fibres (Remark~\ref{rem:vertical}). Equation~\eqref{eq:localrule}
        relates level $n$ to level $n+1$: its shape map is the successor, not the
        identity. It is a horizontal morphism, a different kind of arrow.
  \item \textbf{Two dimensions versus one.} The grid has two directions, and that is
        where the commutator of Remark~\ref{rem:weyl} comes from --- $U$ along one
        axis, $D$ along the other. A bare recurrence is one-dimensional and has no
        commutator to witness.
\end{enumerate}
\end{warning}

The resolution is that the second dimension is available, but only once one leaves
single recurrences for pairs.

\begin{remark}[Wilf--Zeilberger supplies the missing direction]\label{rem:WZ}
The method of \cite{AeqB} proves $\sum_k F(n,k) = \text{const}$ by exhibiting a
certificate $G$ with
\[
  F(n+1,k) - F(n,k) \;=\; G(n,k+1) - G(n,k),
\]
a discrete divergence-free condition on $\NN^2$. Two variables, two directions, and an
identity asserting that two ways round a face agree --- which is the shape of
Remark~\ref{rem:weyl}, not the shape of \eqref{eq:localrule}. Gosper's algorithm,
being univariate, has no analogue here at all; Zeilberger's bivariate extension is
the one that matches.
\end{remark}

\begin{proposition}[\emph{ours} --- the honest form]\label{prop:honest}
A Fomin local rule is a bijection witnessing $DU - UD = I$. A Wilf--Zeilberger pair is
a bijection witnessing a telescoping identity. The common structure is: \emph{a
commutator identity categorified on a two-dimensional growth diagram}. This is a
comparison between Fomin's rules and WZ \emph{pairs}, and not between Fomin's rules
and hypergeometric series as such.
\end{proposition}

\begin{question}\label{q:wz}
Can a WZ certificate be read as an explicit container isomorphism --- not merely as
evidence that a bijective proof exists, but as the bijection? Proposition~\ref{prop:honest}
fixes the form the answer must take: a vertical morphism in $\Cont$ over the face data
of a two-dimensional diagram, as in Remark~\ref{rem:vertical}.

By Warning~\ref{warn:analogy} the calibration case must be laid out as a diagram and
not as a recurrence. The natural choice is $q$-Vandermonde
(Proposition~\ref{prop:qvan}), because its bijection is already known ---
Corollary~\ref{cor:levi} --- so the test is sharp: \emph{does the WZ certificate for
$q$-Vandermonde reproduce the subspace bijection?} If it does, WZ is a machine for
producing container isomorphisms. If it does not, the two notions of ``local rule''
are merely cousins and this section should end at Warning~\ref{warn:analogy}.
\end{question}

\begin{remark}[a second route to the missing trace]\label{rem:trace}
Recorded because it arrives here from a direction independent of the containers work.
Thesis \S4.3.1 states that each application of the local rule corresponds to an
application of the commutation relation \emph{and} that each application of the
rotation operator corresponds to the cylindric invariance of the trace. If the first
correspondence is Remark~\ref{rem:weyl}, the second says the cylindric case requires a
\emph{trace} on the container category --- which is the open
missing-trace-on-$\con$ question from the containers work, reached from
representation theory rather than from category theory. Two independent routes to one
gap is worth more than one.
\end{remark}

\section{At $q$, the notation becomes honest}

\subsection{The obligatory factor is an automorphism group order}

With $(a;q)_n = \prod_{i=0}^{n-1}(1 - aq^i)$, basic hypergeometric series are
\[
  {}_p\phi_{q}\!\left(\begin{matrix}a_1,\dots,a_p\\ b_1,\dots,b_q\end{matrix};q,z\right)
  = \sum_{n\ge0}
  \frac{(a_1;q)_n\cdots(a_p;q)_n}{(q;q)_n\,(b_1;q)_n\cdots(b_q;q)_n}\,z^n ,
\]
and the obligatory factor is now $(q;q)_n$ rather than $n!$ \cite{GasperRahman}.

\begin{proposition}\label{prop:qpoch-GL}
$\displaystyle \norm{\GL_n(\FF_q)} = \prod_{i=0}^{n-1}(q^n - q^i)
 = q^{\binom{n}{2}}\prod_{j=1}^{n}(q^j-1)
 = (-1)^n q^{\binom{n}{2}}(q;q)_n .$
\end{proposition}

\begin{corollary}[\emph{ours}]\label{cor:qhonest}
The obligatory $(q;q)_n$ of ${}_p\phi_q$ is $\norm{\GL_n(\FF_q)}$ up to a sign and the
factor $q^{\binom{n}{2}}$. It is therefore the order of the automorphism group of the
$n$-dimensional $q$-container, and dividing by it is a genuine quotient --- case (i)
of Proposition~\ref{prop:downstairs}, not case (ii).
\end{corollary}

So feature (c), which is only accidentally structural at $q = 1$ (where it happens
that $(1)_n = \norm{S_n}$), is structural at $q$ \emph{by construction}. And the
discrepancy is exactly the familiar one:

\begin{remark}[the torus again]\label{rem:torus}
$q^{\binom{n}{2}}$ is the unipotent part and the remaining $(q-1)^n$ inside
$(q;q)_n = \prod(1-q^j)$ is the torus. \texttt{q-dirichlet.tex} identifies the torus
factor $(q-1)^{\dim}$ as precisely what makes the additive leg of Day convolution work
and the multiplicative leg fail. It reappears here as the discrepancy between the
$q$-Pochhammer and the group order, which is a reason to expect that the obstruction
of that note and the normalisation of basic hypergeometric series are the same
phenomenon seen twice.
\end{remark}

\noindent
Note also that $\norm{\GL_n(\FF_q)}$ is the divisor in the $q$-Dirichlet series
$\sum_{n\ge1}a_n n^{-s}/\norm{\GL_n(\FF_q)}$ of \texttt{q-dirichlet.tex}. The
``spare $n!$'' of Definition~\ref{def:pFq} and that divisor are the same object at
$q$.

\subsection{$q$-Vandermonde is already in the table}

Gaussian binomials are honest $q$-containers: $\qbin{n}{k}$ counts $k$-dimensional
subspaces of $\FF_q^n$, so it is a shape-count, not a quotient, and
\[
  \qbin{n}{k} = \frac{(q;q)_n}{(q;q)_k\,(q;q)_{n-k}} .
\]

\begin{proposition}[$q$-Vandermonde]\label{prop:qvan}
$\displaystyle \qbin{m+n}{r} = \sum_{k}\qbin{m}{k}\qbin{n}{r-k}\,q^{(m-k)(r-k)} .$
\end{proposition}

\begin{example}\label{ex:qvan}
$m = n = r = 1$: left side $\qbin{2}{1} = 1+q$; right side
$k=0$ gives $1\cdot1\cdot q^{1} = q$ and $k=1$ gives $1\cdot1\cdot q^{0} = 1$. \OK
\end{example}

\begin{corollary}[\emph{ours}]\label{cor:levi}
Proposition~\ref{prop:qvan} is the Levi row of the decategorification table of
\texttt{q-container.tex}: ``Day convolution along $\oplus$ (Levi) $\leftrightarrow$
multiplication weighted by $q^{k(n-k)}$''. Its combinatorial content is the
classification of an $r$-subspace of $\FF_q^{m+n}$ by its intersection pattern with a
fixed $m$-dimensional summand, which is a container isomorphism over the dimension
map. So the first non-trivial summation identity of the basic hypergeometric canon is
a theorem already proved at the container level in that note.
\end{corollary}

\begin{corollary}[the thesis of \S6]\label{cor:thesis}
Basic hypergeometric series are more container-native than classical ones. At $q = 1$
the denominator is a quotient only in the single case $b = 1$; at $q$ it is an
automorphism group order by construction, and the structure constants of the series
are subspace counts. The $q$-deformation is not a complication of the classical
theory here --- it is the setting in which the classical theory's conventions become
statements.
\end{corollary}

\section{Where it stops}

Everything above needs $a_i, b_j \in \NN$. For
${}_2F_1(\tfrac12,-\tfrac13;\pi;z)$ there is no container anywhere near it, and its
coefficients are not even rational.

\begin{remark}[\emph{ours} --- $\Gamma$ interpolates past the containers]\label{rem:gamma-past}
This is not a defect to be patched. \texttt{gamma.tex} \S1 characterises $\Gamma$ as
the unique log-convex solution of $f(x+1) = xf(x)$ --- the smooth curve through the
factorials. Transposed into the present language: hypergeometric parameters at
non-negative integers are container data, and $\Gamma$ is what carries the apparatus
to parameters at which no container exists. The analytic theory of ${}_pF_q$ is the
analytic continuation, in its parameters, of a combinatorial theory. That is why the
subject looks as though it came from nowhere, and it is the same boundary as the two
obstructions of \texttt{list-and-bag.tex}: there the \emph{functional} $\int_\RR$ was
not a container operation, here the \emph{parameters} are not container data.
\end{remark}

\begin{warning}\label{warn:preserve}
The disagreement should be preserved rather than resolved. There is no useful sense in
which a general ${}_pF_q$ is a container; there is a precise sense in which the
integer-parameter ones are; and $\Gamma$ is the interpolation between. Claiming more
than this --- that hypergeometric series ``are'' containers --- would be the error the
rest of \texttt{q-bool/} is careful to avoid.
\end{warning}

\section{Next}

\begin{enumerate}[nosep]
  \item \textbf{Question~\ref{q:had-q}: the Hadamard product at $q$.} Cheapest item
        here. Pull back along $\dim$, check closure, check that it decategorifies to
        the coefficientwise product of $q$-series. If it does, the table of
        \texttt{q-container.tex} gains a row and Warning~\ref{warn:column} gains its
        $q$-analogue.
  \item \textbf{Add the Hadamard column} to the table of \texttt{q-dirichlet.tex}
        \S2, with Warning~\ref{warn:column}: the three rows there are Day
        convolutions, and this is the structure that is not one.
  \item \textbf{Question~\ref{q:wz}: WZ certificates as container isomorphisms.}
        The highest-value item, and the one that connects to \texttt{lascoux/}. Start
        with a single identity whose certificate is integral --- the $q$-Vandermonde
        of Proposition~\ref{prop:qvan} has a known bijective proof, so it is the
        calibration case: does the WZ certificate for it reproduce the subspace
        bijection of Corollary~\ref{cor:levi}?
  \item \textbf{Refute or confirm Conjecture~\ref{conj:no2F1}} by computing the
        Mellin transform of ${}_2F_1(a,b;c;-x^2)$ and inspecting the
        $c$-dependence. This bounds the programme and should be done before it is
        extended.
\end{enumerate}

\begin{thebibliography}{9}

\bibitem{AeqB} Marko Petkov\v{s}ek, Herbert S.\ Wilf and Doron Zeilberger,
\emph{$A=B$}, A.\ K.\ Peters, 1996. Gosper's algorithm, Zeilberger's algorithm, and
the WZ method.

\bibitem{BLL} Fran\c{c}ois Bergeron, Gilbert Labelle and Pierre Leroux,
\emph{Combinatorial Species and Tree-like Structures}, Cambridge University Press,
1998. The Cartesian product of species is \S2.1.

\bibitem{Day} Brian Day, \emph{On closed categories of functors}, Reports of the
Midwest Category Seminar IV, Lecture Notes in Math.\ 137, Springer, 1970.

\bibitem{GasperRahman} George Gasper and Mizan Rahman, \emph{Basic Hypergeometric
Series}, 2nd ed., Cambridge University Press, 2004.

\bibitem{Joyal81} Andr\'e Joyal, \emph{Une th\'eorie combinatoire des s\'eries
formelles}, Adv.\ Math.\ \textbf{42} (1981), 1--82.

\bibitem{Kock} Joachim Kock, \emph{Data types with symmetries and polynomial functors
over groupoids}, ENTCS \textbf{286} (2012), 351--365.

\bibitem{Morrison} Kent E.\ Morrison, \emph{An introduction to $q$-species},
Electron.\ J.\ Combin.\ \textbf{12} (2005), R62.

\end{thebibliography}

\end{document}
